\documentclass[a4paper,10pt]{extarticle}
\usepackage[T2A]{fontenc}
\usepackage[utf8]{inputenc}
\usepackage[english,russian]{babel}
\usepackage{paratype}
% Достаточно для нормальной печати
\usepackage[margin=15mm]{geometry}
\usepackage{xcolor}
\usepackage{tikz}
\pagestyle{empty}

% Название олимпиады и тура
\newcommand{\OlympiadName}{Всероссийская олимпиада школьников по физике}
\newcommand{\OlympiadDetails}{Региональный этап}
% Ширина текстовых блоков внутри рамки заголовка
\newcommand{\ContentWidth}{115mm}
% Геометрия рамки заголовочного блока (название, чекбоксы, инструкция) --
% фиксированные координаты нужны и для рисования, и для последующего
% алгоритмического поиска этого блока на скане
\newcommand{\FrameLeft}{15}
\newcommand{\FrameRight}{140}
\newcommand{\FrameTop}{7}
\newcommand{\FrameBottom}{63}
\newcommand{\FramePad}{3}
% Угловая реперная метка -- жирный штрих ПО САМОЙ рамке (часть её контура),
% а не отдельная фигура сбоку. #1,#2 -- угол рамки, #3,#4 -- знаки направления
% вдоль рамки от угла (+1/-1), т.е. куда "растут" обе стороны уголка.
\newcommand{\CornerMark}[4]{%
  \draw[fill=black,draw=none]
    (#1-1.1*#3,#2-1.1*#4) -- (#1+12*#3,#2-1.1*#4) --
    (#1+12*#3,#2+1.1*#4) -- (#1+1.1*#3,#2+1.1*#4) --
    (#1+1.1*#3,#2+12*#4) -- (#1-1.1*#3,#2+12*#4) -- cycle;
}

\newcommand{\Instructions}{
  На одной странице можно оформлять \textbf{только одну задачу}.
  Отметьте \textbf{крестиком} номер задачи решение которой представлено на странице. Также укажите номер и суммарное количество страниц в \textbf{решении каждой задачи отдельно.}

  Если отмечено более одной задачи или не отмечено ни одной задачи, страница считается черновиком и не проверяется.
  \textbf{Страницы по каждой из задач нумеруются отдельно.}
}

\begin{document}
\null
\begin{tikzpicture}[remember picture,overlay]
  \begin{scope}[shift={(current page.north west)},x=1mm,y=-1mm]
    % Мероприятие и тур
    \node[anchor=north west,inner sep=0pt,text width=\ContentWidth,
          font=\fontsize{11}{13}\selectfont\bfseries]
      at (\FrameLeft+\FramePad,10) {\mbox{\OlympiadName}};
    \node[anchor=north west,inner sep=0pt,text width=\ContentWidth,
          font=\fontsize{11}{13}\selectfont]
      at (\FrameLeft+\FramePad,15) {\OlympiadDetails};

    \foreach \i in {1,2,...,5} {
        \node[font=\fontsize{16}{18}\selectfont] at (12+12*\i,24) {$\i$};
        \draw (9+12*\i,29) rectangle ++(6,6);
    }

    % "Страница __ из __" -- в одной строке с рядом чекбоксов отметки задач
    \node[anchor=west,inner sep=0pt,font=\fontsize{12}{14}\selectfont]
      at (85,32) {Страница~\underline{\hspace{7mm}}~из~\underline{\hspace{7mm}}};

    % Угловые реперные метки по углам рамки заголовочного блока --
    % для автоматической привязки (поиска) блока на скане
    \CornerMark{\FrameLeft}{\FrameTop}{1}{1}
    \CornerMark{\FrameRight}{\FrameTop}{-1}{1}
    \CornerMark{\FrameLeft}{\FrameBottom}{1}{-1}
    \CornerMark{\FrameRight}{\FrameBottom}{-1}{-1}

    % Место под QR-код/шифр: 48 x 40 мм
    \draw[dash pattern=on 2mm off 1mm,line width=0.5pt] (147,12) rectangle ++(48,40);
    \draw[black!30,line width=0.5pt] (151,32) -- (191,32);
    \node[anchor=north,inner sep=0pt,text width=42mm,align=center,text=black!30,
          font=\fontsize{9}{11}\selectfont]
      at (171,33.5) {Шифр\\[0.5mm]\fontsize{6.5}{8}\selectfont(заполняется шифровальной\\ комиссией)};

    % Рамка вокруг заголовочного блока (название, чекбоксы, инструкция) --
    % обязательный элемент макета для алгоритмического поиска на скане
    \draw[line width=1pt,black!85] (\FrameLeft,\FrameTop) rectangle (\FrameRight,\FrameBottom);

    \node[anchor=north west,inner sep=0pt,text width=\ContentWidth,
          font=\fontsize{8.5}{10.2}\selectfont]
      at (\FrameLeft+\FramePad,38) {\Instructions};

    % Миллиметровка: 180 x 220 мм (та же область, что и у клетчатого бланка)
    \begin{scope}[shift={(15,68)}]
      \foreach \x in {1,2,...,179}
        \draw[line width=0.6pt,black!15] (\x,0) -- (\x,220);
      \foreach \y in {1,2,...,219}
        \draw[line width=0.6pt,black!15] (0,\y) -- (180,\y);
      \foreach \x in {5,10,...,175}
        \draw[line width=0.75pt,black!22] (\x,0) -- (\x,220);
      \foreach \y in {5,10,...,215}
        \draw[line width=0.75pt,black!22] (0,\y) -- (180,\y);
      \foreach \x in {10,20,...,170}
        \draw[line width=1pt,black!32] (\x,0) -- (\x,220);
      \foreach \y in {10,20,...,210}
        \draw[line width=1pt,black!32] (0,\y) -- (180,\y);
      \foreach \x in {50,100,150}
        \draw[line width=1.25pt,black!45] (\x,0) -- (\x,220);
      \foreach \y in {50,100,150,200}
        \draw[line width=1.25pt,black!45] (0,\y) -- (180,\y);

      % Граница
      \draw[line width=1.5pt,black!55] (0,0) rectangle (180,220);
    \end{scope}
  \end{scope}
\end{tikzpicture}
\end{document}
